% Source: https://artofproblemsolving.com/community/c3h2178294_vietas_formula
% Stable identity: aops-topic-2178294-q01
Since $a\ne2$, the constant term $2-a$ is nonzero, so $0$ is not a solution of the given equation. If $w$ is a solution, divide $(2+a)w^2+5w+(2-a)=0$ by $w^2$:
\[(2-a)\left(\dfrac{1}{w}\right)^2+5\cdot\dfrac{1}{w}+(2+a)=0 .\]
So the reciprocals of the solutions are the solutions of $(2-a)x^2+5x+(2+a)=0$, a quadratic since $a\ne2$. Two quadratics with the same roots (counted with multiplicity) are constant multiples of each other, so
\[(r,s,t)=k\left(2-a,\ 5,\ 2+a\right)\]
for some constant $k$. A common divisor of $2-a$, $5$ and $2+a$ divides $5$ and $(2-a)+(2+a)=4$, so these three integers have greatest common factor $1$. Writing $k=\dfrac{m}{q}$ in lowest terms, $q$ must divide all three, so $q=1$; then the greatest common factor of $r,s,t$ is $|k|$, so $k=\pm1$, and $s=5k>0$ gives $k=1$. Hence
\[r+s+t=(2-a)+5+(2+a)=\boxed{9},\]
whatever the value of $a$.
